\section{Conclusion}

This review systematized 41 studies on AI-assisted software development frameworks that shift the unit of work from the isolated prompt to a structured process, with persistent artifacts, agentic roles, integrated execution, and validation. The thesis that this shift has already happened rests on the evidence: the contribution is described as a framework in 37 of the 41 studies, and the declared differentiator is the process structure, not the underlying language model. On that corpus, we proposed and tested a six-dimension taxonomy. Five dimensions (specification, context, roles, execution, and validation) receive strong empirical support; portability is weak and predominantly diagnostic, and two relevant axes, the human-organizational and the security one, emerge outside the taxonomy and recommend extending it.

The field is young, has accelerated since 2024, and still lacks a consolidated publication forum, which is reflected in the dispersion of venues and the sizable presence of preprints. The evidence is uneven, and three axes of tension (autonomy versus human control, heavy governance versus lightweight agility, and specifying first versus recovering the specification from legacy) structure the debate without resolving it. The most cited practical bottlenecks are not in code generation but in keeping the agent anchored in the real project and auditable: context and traceability are the dominant transversal requirements.

The consequence for future research is direct. The field's next step is empirical and process-oriented: benchmarks that evaluate the pipeline and not only the solution, direct comparisons between frameworks, portability instruments, longitudinal evaluation of governance, a sociotechnical layer in the classification, and systematic security evaluation. This review offers the validated taxonomy and the agenda that organize that effort.
